% Source-derived chapter 03: Strong perversity for Chern classes
% Original source: pw-conjecture-gl-n
\chapter{Strong perversity for Chern classes}
\label{pw-conjecture-gl-n-chapter-03}
\source{pw-conjecture-gl-n}

\begin{remark}[Source section]
\label{pw-conjecture-gl-n-chapter-03-source}
\source{pw-conjecture-gl-n}
\source{pw-conjecture-gl-n}
This chapter reproduces the corresponding section of the retrieved
LaTeX source, including its definitions, statements, examples, and
proofs. It has not yet been translated into Lean.
\notready
\end{remark}

% The source section title was: Strong perversity for Chern classes
In this section, we fix the rank $n$ and the degree $d$ with $(n,d)=1$. In Theorem \ref{conj2.7}, we rephrase and then enhance Conjecture \ref{conj} to a statement involving $\CL$-twisted Hitchin systems and strong perversity of Chern classes. It will be proven in Sections 3 and 4.

\subsection{Tautological classes}

As discussed in \cite[Section 0.3]{dCMS}, the $P=W$ conjecture for $\mathrm{GL}_n$ can be reduced to a statement involving tautological classes on $M_{\mathrm{Dol}}$ and the perverse filtration associated with $h: M_{\mathrm{Dol}} \to A$, without reference to the Betti moduli space $M_B$. In this subsection, we recall this reduction step.


For convenience, we work with the $\mathrm{PGL}_n$ Dolbeault moduli space to avoid normalization of a universal family as in \cite{dCMS}; we refer to \cite{Shende} for a detailed discussion concerning the formulation of the $P=W$ conjecture in terms of tautological classes for the $\mathrm{PGL}_n$ Dolbeault moduli space.

Fix $\CN \in \mathrm{Pic}^d(C)$. Let $\widecheck{M}_{\mathrm{Dol}}$ be 
the moduli stack of stable Higgs bundles $(\CE, \theta)$ with $\mathrm{det}(\CE) \simeq \CN$ and $\mathrm{trace}(\theta) = 0$, rigidified with respect to the generic $\mu_n$-stabilizer; this is the same as taking its coarse moduli space.  We refer to this (nonsingular) variety as the $\mathrm{SL}_n$ Dolbeault moduli space of degree $d$. The finite group $\Gamma: = \mathrm{Pic}^0(C)[n]$ acts naturally on $\widecheck{M}_{\mathrm{Dol}}$ via tensor product. The $\mathrm{PGL}_n$ Dolbeault moduli space of degree $d$ is recovered by taking the quotient stack
\begin{equation}\label{233}
\source{pw-conjecture-gl-n}
\widehat{M}_{\mathrm{Dol}} : = \widecheck{M}_{\mathrm{Dol}}/\Gamma
\end{equation}
which is a nonsingular Deligne--Mumford stack. The $\mathrm{PGL}_n$ Hitchin system 
\[
\widehat{h}: \widehat{M}_{\mathrm{Dol}} \to \widehat{A}: = \bigoplus_{i=2}^n H^0(C, {\Omega_C^1}^{\otimes i})
\]
is induced by the Hitchin map associated with $\widecheck{M}_{\mathrm{Dol}}$ as the $\Gamma$-action is fiberwise with respect to $h$. Analogous to the $\mathrm{GL}_n$  case, we have the perverse filtration $P_*H^*(\widehat{M}_{\mathrm{Dol}}, \BQ)$ associated with $\widehat{h}$. The universal $\mathrm{PGL}_n$-bundle $\CU$ on $C \times \widehat{M}_{\mathrm{Dol}}$ induces Chern characters
\[
\mathrm{ch}_k(\CU) \in H^{2k}(C \times \widehat{M}_{\mathrm{Dol}}, \BQ),\quad k \geq 2. 
\]
The tautological classes $c_k(\gamma)$ are defined to be
\[
c_k(\gamma): = \int_\gamma \mathrm{ch}_k(\CU) = q_{M*}(q_C^*\gamma \cup \mathrm{ch}_k(\CU)) \in H^*(\widehat{M}_{\mathrm{Dol}}, \BQ), \quad \gamma \in H^*(C,\BQ),
\]
where $q_{(-)}$ are the projections from $C \times \widehat{M}_{\mathrm{Dol}}$.

Now we consider the $\mathrm{PGL}_n$ Betti moduli space of degree $d$ with the isomorphism on cohomology provided by non-abelian Hodge theory (\emph{c.f.}\cite[Theorem 1.2.4]{dCHM1}):
\begin{equation}\label{NAH2}
\source{pw-conjecture-gl-n}
H^*(\widehat{M}_{\mathrm{Dol}}, \BQ) = H^*(\widehat{M}_{B}, \BQ).
\end{equation}
Define the Hodge sub-vector space
\[
^k\mathrm{Hdg}^m(\widehat{M}_B): = W_{2k}H^m(M_B, \BQ) \cap F^kH^m(M_B, \BC) \subset H^m(\widehat{M}_{B}, \BQ).
\]

The following theorem, collecting results of Markman and Shende, provides a complete description of $H^*(\widehat{M}_B, \BQ)$ in terms of the Chern classes $\mathrm{ch}_k(\CU)$ and the weight filtration.

\begin{thm}[\cite{Markman, Shende}]
\source{pw-conjecture-gl-n}
\notready\label{thm2.1}
\source{pw-conjecture-gl-n}
We use the same notation as above.
\begin{enumerate}
    \item[(i)] The tautological classes $c_k(\gamma)\in H^*(\widehat{M}_{\mathrm{Dol}}, \BQ)$ generate $H^*(\widehat{M}_\mathrm{Dol}, \BQ)$ as a $\BQ$-algebra.
    \item[(ii)] The class $c_k(\gamma)$, passing through the non-abelian Hodge correspondence (\ref{NAH2}), lies in $^k\mathrm{Hdg}^*(\widehat{M}_B)$. In particular, we have a canonical decomposition
    \[
H^*(\widehat{M}_B, \BQ) = \bigoplus_{m,k} {^k\mathrm{Hdg}^m(\widehat{M}_B)}.    \]
\end{enumerate}
\end{thm}

\begin{proof}
The first part was proven in \cite{Markman}, and the second part was proven in \cite{Shende}.
\end{proof}


Theorem \ref{thm2.1} (ii) yields immediately that 
\[
W_{2k}H^m(\widehat{M}_B, \BQ) = W_{2k+1}H^m(\widehat{M}_B, \BQ).
\]
Moreover, by Theorem \ref{thm2.1}, the $P=W$ conjecture implies that each class $\prod_{i=1}^s c_{k_i}(\gamma_i)$ lies in the perverse piece $P_{\Sigma_i{k_i}}H^*(\widehat{M}_{\mathrm{Dol}}, \BQ)$; the latter is in fact equivalent to the $P=W$ conjecture. Indeed, suppose we have 
\begin{equation}\label{taut}
\source{pw-conjecture-gl-n}
\prod_{i=1}^s c_{k_i}(\gamma_i) \in P_{\Sigma_i{k_i}}H^*(\widehat{M}_{\mathrm{Dol}}, \BQ)
\end{equation}
for any product of tautological classes.
Then we know that $W_{2k}H^*(\widehat{M}_B, \BQ) \subset P_kH^*(\widehat{M}_{\mathrm{Dol}}, \BQ)$. The curious hard Lefschetz theorem proven by Mellit \cite{Mellit} forces the two filtrations $P_\bullet$ and $W_{2\bullet}$ to coincide as long as one contains the other. This was mentioned in the last paragraph of \cite[Section 1]{Mellit}; we include its proof here for the reader's convenience.

\begin{lem}
\source{pw-conjecture-gl-n}
\notready
We denote by $V$ the $\BQ$-vector space (\ref{NAH2}) with the perverse and the weight filtrations $P_\bullet$ and $W_{2\bullet}$. If $W_{2k} \subset P_k$ for all $k$, then $W_{2k} = P_k$ for all $k$.
\end{lem}

\begin{proof}
Assume $r = \mathrm{dim}\widehat{M}$; both filtrations terminate at the $r$-th pieces, \emph{i.e.}, $W_{2r} = P_{r} = V$. We first show that 
\[
W_0 = P_0, \quad W_{2(r-1)} = P_{r-1}. 
\]
In fact, since $W_\bullet  \subset P_\bullet$, we have
\[
\mathrm{dim} W_0 \leq \mathrm{dim} P_0 = \mathrm{dim}V/P_{r-1} \leq \mathrm{dim} V/W_{2(r-1)};
\]
moreover, by the curious hard Lefschetz theorem, each inequality has to be an equality. So our claim follows. 

Then we proceed by applying the same argument to $W_1, P_1$ and $W_{2(r-2)}, P_{r-1}$. The lemma follows by a simple induction.
\end{proof}


In conclusion, we have reduced the $P=W$ conjecture to the following:


\begin{conj}[Equivalent version of P=W]
\source{pw-conjecture-gl-n}
\notready\label{conj2.2}
\source{pw-conjecture-gl-n}
Condition (\ref{taut}) holds for all products of tautological classes.
\end{conj}

\subsection{Strong perversity for Chern classes for $\CL$-twisted Hitchin systems}

For our purposes, it is important to consider Dolbeault moduli spaces of Higgs bundles, twisted by an effective line bundle $\CL$ (\emph{i.e.}, $H^0(C, \CL) \neq 0$). These moduli spaces have already appeared in \cite{CL, HT2, HT, Yun1, Yun2}; we review the construction here briefly. 

Set $\Omega_\CL$ to be the line bundle $\Omega^1_C \otimes \CL$ on the curve $C$. We denote by $M^\CL_{\mathrm{Dol}}$ the moduli space of stable twisted Higgs bundles
\[
(\CE, \theta), \quad \theta: \CE \to \CE \otimes\Omega_\CL, \quad \mathrm{rk}(\CE) = n, ~~\mathrm{deg}(\CE)=d,
\]
with respect to the slope stability condition. The corresponding Hitchin map
\[
h^\CL: M^\CL_{\mathrm{Dol}} \to A^\CL: = \bigoplus_{i=1}^n H^0\left(C, {\Omega_\CL}^{\otimes i} \right) ,\quad (\CE, \theta) \mapsto \mathrm{char.polynomial}(\theta),
\]
is still proper as in the untwisted case; but it fails to be a Lagrangian fibration when $\mathrm{deg}(\CL)>0$. The $\CL$-twisted $\mathrm{SL}_n$ and $\mathrm{PGL}_n$ Dolbeault moduli spaces $\widecheck{M}^\CL_{\mathrm{Dol}}$ and $\widehat{M}^\CL_{\mathrm{Dol}}$ can be constructed similarly. The moduli space
\[
\widecheck{M}^\CL_{\mathrm{Dol}}:= \{(\CE, \theta) \in M^\CL_{\mathrm{Dol}}|~~ \mathrm{det}(\CE) \simeq \CN \in \mathrm{Pic}^d(C), \mathrm{trace}(\theta) = 0\}
\]
admits a Hitchin map
\[
\widecheck{h}^\CL: \widecheck{M}^\CL_{\mathrm{Dol}} \to \widehat{A}^\CL:= \bigoplus_{i=2}^nH^0\left(C, {\Omega_\CL}^{\otimes i} \right)\]
and a fiberwise $\Gamma = \mathrm{Pic}^0(C)[n]$ action by tensor product. Taking the $\Gamma$-quotient recovers the $\CL$-twisted $\mathrm{PGL}_n$ Hitchin map
\[
\widehat{h}^\CL: \widehat{M}^\CL_{\mathrm{Dol}} =\widecheck{M}^\CL_{\mathrm{Dol}}/\Gamma \to \widehat{A}^\CL. 
\]


An observation in \cite[Section 4]{MS} is that, for a fixed closed point $p\in C$, the $\CL$-twisted and $\CL(p)$-twisted $\mathrm{SL}_n$ Dolbeault moduli spaces can be related via critical loci and vanishing cycles, which we recall in the following.

By viewing a $\CL$-twisted Higgs bundle naturally as a $\CL(p)$-twisted Higgs bundle, we have the natural embedding $i: \widecheck{M}^\CL_{\mathrm{Dol}} \hookrightarrow \widecheck{M}^{\CL(p)}_{\mathrm{Dol}}$ which induces the commutative diagram
\begin{equation}\label{2.3_1}
\source{pw-conjecture-gl-n}
\begin{tikzcd}
\widecheck{M}^\CL_{\mathrm{Dol}}\arrow[r, hook,"i"] \arrow[d, "\widecheck{h}^\CL"]
&\widecheck{M}^{\CL(p)}_{\mathrm{Dol}} \arrow[d, "\widecheck{h}^{\CL(p)}"] \\
\widehat{A}^\CL \arrow[r,hook]
& \widehat{A}^{\CL(p)}.
\end{tikzcd}
\end{equation}

We recall the following theorem from \cite{MS}.

\begin{thm}[\cite{MS} Theorem 4.5]
\source{pw-conjecture-gl-n}
\notready\label{thm2.5}
\source{pw-conjecture-gl-n}
There exists a regular function $g: \widecheck{M}^{\CL(p)}_{\mathrm{Dol}} \to \BA^1$ factorized as $g = \nu \circ \widecheck{h}^{\CL(p)}$ with $\nu: \widehat{A}^{\CL(p)} \to \BA^1$ such that
\[
\varphi_g \simeq \mathrm{IC}_{\widecheck{M}^{\CL}_{\mathrm{Dol}}} = \BQ_{\widecheck{M}^{\CL}_{\mathrm{Dol}}}[\mathrm{dim} \widecheck{M}^{\CL}_{\mathrm{Dol}}].
\]
\end{thm}

Since the embedding $i: \widecheck{M}^\CL_{\mathrm{Dol}} \hookrightarrow \widecheck{M}^{\CL(p)}_{\mathrm{Dol}}$ is $\Gamma$-equivariant, taking $\Gamma$-quotients in the diagram (\ref{2.3_1}) yields the following commutative diagram:
\begin{equation}\label{diagram3}
\source{pw-conjecture-gl-n}
\begin{tikzcd}
\widehat{M}^\CL_{\mathrm{Dol}}\arrow[r, hook, "\widehat{i}"] \arrow[d, "\widehat{h}^\CL"]
&\widehat{M}^{\CL(p)}_{\mathrm{Dol}} \arrow[d, "\widehat{h}^{\CL(p)}"] \\
\widehat{A}^\CL \arrow[r, hook]
& \widehat{A}^{\CL(p)}.
\end{tikzcd}
\end{equation}


\begin{prop}
\source{pw-conjecture-gl-n}
\notready\label{prop2.6}
\source{pw-conjecture-gl-n}
The vanishing cycle complex $\varphi_{\widehat{g}}$ associated with the regular function
\[
\widehat{g}: = \nu \circ \widehat{h}^{\CL(p)}: \widehat{M}^{\CL(p)}_{\mathrm{Dol}} \to \BA^1
\]
satisfies
\[
\varphi_{\widehat{g}} \simeq \mathrm{IC}_{\widehat{M}^{\CL}_{\mathrm{Dol}}} = \BQ_{\widehat{M}^{\CL}_{\mathrm{Dol}}}\left[\mathrm{dim} \widehat{M}^{\CL}_{\mathrm{Dol}}\right].\]
\end{prop}

\begin{proof}
We reduce Proposition \ref{prop2.6} to Theorem \ref{thm2.5}. Consider the $\Gamma$-quotient map $r: \widecheck{M}^{\CL(p)}_{\mathrm{Dol}} \to \widehat{M}^{\CL(p)}_{\mathrm{Dol}}$. The direct image $r_* \BQ_{\widecheck{M}^{\CL(p)}_{\mathrm{Dol}}}$ admits a natural $\Gamma$-equivariant structure, whose invariant part recovers 
\begin{equation}\label{2.6_1}
\source{pw-conjecture-gl-n}
\left(r_* \BQ_{\widecheck{M}^{\CL(p)}_{\mathrm{Dol}}}\right)^{\Gamma} \simeq \BQ_{\widehat{M}^{\CL(p)}_{\mathrm{Dol}}}.
\end{equation}
Since $g = \widehat{g}\circ r$, we have
\begin{align*}
\varphi_{\widehat{g}} & = \varphi_{\widehat{g}}\left(\BQ_{\widehat{M}^{\CL(p)}_{\mathrm{Dol}}}\left[\mathrm{dim} \widehat{M}^{\CL(p)}_{\mathrm{Dol}}\right]\right) \\ & \simeq \varphi_{\widehat{g}}\left(r_* \BQ_{\widecheck{M}^{\CL(p)}_{\mathrm{Dol}}}\left[\mathrm{dim} \widecheck{M}^{\CL(p)}_{\mathrm{Dol}}\right]\right)^{\Gamma}\\
& \simeq (r_* \varphi_g)^\Gamma \\ & \simeq  \BQ_{\widehat{M}^{\CL}_{\mathrm{Dol}}}\left[\mathrm{dim} \widehat{M}^{\CL}_{\mathrm{Dol}}\right].
\end{align*}
Here the first equation follows by definition, the second uses (\ref{2.6_1}), the third follows from the base change, and the last is given by Theorem \ref{thm2.5}.
\end{proof}

%In the remaining, we show that the following conjecture implies Conjecture \ref{conj2.3}; therefore it also implies Conjecture \ref{conj} in view of Proposition \ref{prop2.41}.

Now we formulate a sheaf-theoretic enhancement of Conjecture \ref{conj2.2}.

\begin{thm}
\source{pw-conjecture-gl-n}
\notready\label{conj2.7}
\source{pw-conjecture-gl-n}
There exists an effective line bundle $\CL$ such that the class
\[
\mathrm{ch}_k(\CU^\CL) \in H^{2k}(C\times \widehat{M}^\CL_{\mathrm{Dol}}, \BQ)
\]
has strong perversity $k$ with respect to 
\[
\mathbf{h}^\CL: = \mathrm{id} \times \widehat{h}^\CL:C\times \widehat{M}^\CL_{\mathrm{Dol}} \to C\times \widehat{A}^\CL. 
\]
Here $\CU^\CL$ is the universal $\mathrm{PGL}_n$-bundle over $C\times \widehat{M}^\CL_{\mathrm{Dol}}$.
\end{thm}

Notice that the perversity bound here is stronger than the trivial one of $2k$.  We prove in the rest of this section that Theorem \ref{conj2.7} indeed implies Conjecture \ref{conj2.2}; therefore it implies Conjecture \ref{conj}. We prove Theorem \ref{conj2.7} in Sections 3 and 4.


\subsection{Theorem \ref{conj2.7} implies Conjecture \ref{conj2.2}}\label{Sec2.3} We start with the following claim.
\source{pw-conjecture-gl-n}

%\begin{prop}
%Conjecture \ref{conj2.7} implies Conjecture \ref{conj2.3}.
%\end{prop}

%\begin{proof}
\medskip
\noindent {\bf Claim.} For a point $p\in C$, if the class
\[
\mathrm{ch}_k(\CU^{\CL(p)}) \in H^{2k}(C\times \widehat{M}^{\CL(p)}_{\mathrm{Dol}}, \BQ) 
\]
has strong perversity $k$ with respect to $\mathbf{h}^{\CL(p)}: C\times \widehat{M}^{\CL(p)}_{\mathrm{Dol}} \to C\times \widehat{A}^{\CL(p)}$, then the class
\[
\mathrm{ch}_k(\CU^{\CL}) \in H^{2k}(C\times \widehat{M}^{\CL}_{\mathrm{Dol}}, \BQ)
\]
has strong perversity $k$ with respect to $\mathbf{h}^{\CL}: C\times \widehat{M}^\CL_{\mathrm{Dol}} \to C\times \widehat{A}^\CL$.

\begin{proof}[Proof of Claim]

We consider the commutative diagram obtained from (\ref{diagram3}):

\begin{equation}\label{diagram4}
\source{pw-conjecture-gl-n}
\begin{tikzcd}
C \times \widehat{M}^\CL_{\mathrm{Dol}}\arrow[r,hook, "\mathbf{i}"] \arrow[d, "\mathbf{h}^\CL"]
&C\times \widehat{M}^{\CL(p)}_{\mathrm{Dol}} \arrow[d, "\mathbf{h}^{\CL(p)}"] \\
C\times \widehat{A}^\CL \arrow[r,hook]
& C\times \widehat{A}^{\CL(p)}.
\end{tikzcd}
\end{equation}
Here $\mathbf{i} = \mathrm{id} \times \widehat{i}$. By definition,  we have $\mathbf{i}^* \CU^{\CL(p)} = \CU^{\CL}$. Therefore 
\[
\mathbf{i}^* \mathrm{ch}_k(\CU^{\CL(p)}) = \mathrm{ch}_k(\CU^\CL).
\]
We know from Proposition \ref{prop2.6} that the vanishing cycle complex associated with the regular function
\[
\mathbf{g}: =  \widehat{g} \circ \mathrm{pr}_M: C\times \widehat{M}^{\CL(p)}_{\mathrm{Dol}} \to\widehat{M}^{\CL(p)}_{\mathrm{Dol}}\to \BA^1
\]
satisfies
\[
\varphi_{\mathbf{g}} \simeq \mathrm{IC}_{C\times \widehat{M}^{\CL}_{\mathrm{Dol}}} = {\BQ}_{C\times \widehat{M}^{\CL}_{\mathrm{Dol}}}\left[\mathrm{dim}~ (C\times \widehat{M}^{\CL}_{\mathrm{Dol}})\right].
\]
Our claim then follows from Proposition \ref{prop1.4} (applied to the diagram (\ref{diagram4}), the class $\mathrm{ch}_k(\CU^{\CL(p)})$, and the map $\mathbf{g}$).
\end{proof}

%We prove Conjecture \ref{conj2.7} in the next two sections. Note that although for our purpose, we only need \emph{one} effective divisor $D$ satisfying Conjecture \ref{conj2.7}, our proof actually shows that it is true for \emph{any} effective $\CL$.



%\subsection{Enhancement I: strong perversity}

As a consequence of the claim, the statement of Theorem \ref{conj2.7} holds for $\CL \simeq \CO_C$, \emph{i.e.}, the class 
\[
\mathrm{ch}_k(\CU) \in H^{2k}(C\times \widehat{M}_{\mathrm{Dol}}, \BQ)
\] 
has strong perversity $k$ with respect to 
\[
\mathbf{h}:= \mathrm{id} \times \widehat{h}: C\times \widehat{M}_{\mathrm{Dol}} \to C\times \widehat{A}.
\]


%We first lift Conjecture \ref{conj2.2} to a sheaf-theoretic statement concerning the decomposition theorem for $\widehat{h}: \widehat{M}_{\mathrm{Dol}} \to \widehat{A}$.

Finally, to deduce the cohomological statement, we note that since $\mathbf{h}$ is the identity map restricting to $C$, its induced perverse filtration can be described as:
\begin{equation}\label{2.2_1}
\source{pw-conjecture-gl-n}
P_kH^*(C \times \widehat{M}_{\mathrm{Dol}}, \BQ) = H^*(C, \BQ) \otimes P_kH^*(\widehat{M}_{\mathrm{Dol}}, \BQ).
\end{equation}
Choose a homogeneous basis of $H^*(C, \BQ)$:
\[
\Sigma=\{\sigma_0, \sigma_1, \dots, \sigma_{2g+2}\}
\]
with Poincar\'e-dual basis $\{\sigma^{\vee}_{0}, \dots, \sigma^\vee_{2g+2}\}$.

We may express $\mathrm{ch}_k(\CU)$ as
\begin{equation}\label{2.4_1}
\source{pw-conjecture-gl-n}
\mathrm{ch}_k(\CU) = \sum_{\sigma\in \Sigma} \sigma^\vee \otimes c_k(\sigma) \in H^*(C, \BQ) \otimes H^*(\widehat{M}_{\mathrm{Dol}}, \BQ).
\end{equation}
Since $\mathrm{ch}_k(\CU)$ has strong perversity $k$, Lemma \ref{lem1.1} implies that
\[
\mathrm{ch}_k(\CU) \cup - : P_sH^*(C\times \widehat{M}_{\mathrm{Dol}}, \BQ) \to P_{s+k}H^{*}(C\times \widehat{M}_{\mathrm{Dol}}, \BQ).
\]
Applying this operator to $\sigma \otimes P_sH^*(\widehat{M}_{\mathrm{Dol}}, \BQ) \subset H^*(C\times \widehat{M}_{\mathrm{Dol}}, \BQ)$ with $\sigma \in \Sigma$, we obtain from (\ref{2.2_1}) and (\ref{2.4_1}) that
\[
c_k(\sigma) \cup -: P_sH^*(\widehat{M}_{\mathrm{Dol}}, \BQ) \to P_{s+k}H^*(\widehat{M}_{\mathrm{Dol}}, \BQ). 
\]
In particular, for any class $\gamma \in H^*(C, \BQ)$ we have
\[
c_k(\gamma): P_sH^*(\widehat{M}_{\mathrm{Dol}}, \BQ) \to P_{s+k}H^*(\widehat{M}_{\mathrm{Dol}}, \BQ)\]
which further yields
\[
\prod_ic_{k_1}(\gamma_i) = \left( \prod_ic_{k_1}(\gamma_i) \right) \cup 1 \in P_{\Sigma_ik_i}H^*(\widehat{M}_{\mathrm{Dol}}, \BQ).
 \]
 Here we used $1\in P_0H^0(\widehat{M}_{\mathrm{Dol}}, \BQ)$ in the last equation, which is given by Lemma \ref{lem0}.

This completes the proof that Theorem \ref{conj2.7} implies Conjecture \ref{conj2.2}. \qed
